% !TeX program = xelatex
\documentclass[11pt,a4paper]{article}
\usepackage[margin=23mm]{geometry}
\usepackage{fontspec}
\setmainfont{Latin Modern Roman}
\usepackage{amsmath,amssymb,mathtools}
\usepackage{xcolor,enumitem,booktabs,longtable,array,xurl,hyperref}
\definecolor{muted}{HTML}{53636C}
\hypersetup{colorlinks=true,urlcolor=blue,linkcolor=blue}
\setlength{\parindent}{0pt}
\setlength{\parskip}{5pt}
\newcommand{\R}{\mathbb R}
\newcommand{\E}{\mathbb E}
\newcommand{\N}{\mathbb N}
\newcommand{\ind}{\mathbf 1}
\newcommand{\Var}{\operatorname{Var}}
\newcommand{\Cov}{\operatorname{Cov}}
\newcommand{\supp}{\operatorname{supp}}
\DeclareMathOperator*{\argmin}{arg\,min}
\DeclareMathOperator*{\argmax}{arg\,max}
\newcommand{\entrymeta}[3]{{\sffamily\small\color{muted}\textbf{\texttt{#1}}\quad|\quad #2\par #3\par}}
\newcommand{\Problemref}[1]{\texttt{#1}}
\newcommand{\JELref}[2]{\texttt{#2} (JEL \texttt{#1})}
\begin{document}
\section*{Climate tipping points and nonlinear damages}
\textbf{Problem ID:} \texttt{Q54-223}\quad \textbf{JEL:} \texttt{Q54}\par
% BEGIN ECONOMICS STATEMENT
\label{problem:OP-0112}
\label{atlas:prob:C2}

\noindent\textbf{Original identifier:} C2 (atlas item 112). \textbf{Field:} Environmental and climate economics. \textbf{Classification:} editorial research family with established restricted solutions; current open status not certified.

\subsection*{Definitions and assumptions}

Consider dates $t=0,\ldots,H$, where $H\geq1$ is a supplied finite integer and each period has a supplied physical duration $\Delta>0$. Decisions and flow welfare occur at $t=0,\ldots,H-1$; date $H$ has terminal payoff only. For each member of the research family, supply a measurable climate-economic state space $X$, a measurable action space $B$, and nonempty feasible sets $A(x,j)\subseteq B$. Write $s_t=(x_t,J_t)\in S=X\times\{0,1\}^K$, where $K\geq1$ is a supplied finite integer indexing specified physical elements. Declare state coordinates and units, the temperature variable and reference baseline, and whether $J_{t,k}=1$ records threshold crossing, commitment to a transition, or its completion. Any transition progress, elapsed time or delayed damage needed for the dynamics must be represented in the state or history used by the model. Irreversibility over this horizon means $J_{t+1,k}\ge J_{t,k}$ for every $k$; every admissible transition kernel must assign probability one to $j'\geq j$ componentwise. Reversibility or hysteresis needs another state representation.

A transition kernel $Q(dx',dj'\mid x,j,a)$ contains temperature-dependent hazards and interactions. Hazards are conditional probabilities over the declared period, rather than unspecified instantaneous rates; any rate-to-probability conversion must be supplied. Specify joint transitions, including dependence between physical elements, rather than only marginal hazards. Per-period welfare $U:S\times B\to\R$, restricted to feasible actions, incorporates consumption, ecological services and action costs in a declared utility scale; discount factor $\beta\in(0,1]$ is per period and terminal value $G:S\to\R$ is supplied. A policy $\pi$ is a measurable mapping from the observed history into feasible actions. Each model must declare observations and their timing before an action, an initial state or distribution, and its physical-evidence compatibility criterion and economic propagation restrictions. These declarations are model inputs; the cited sources do not uniquely supply them for this research family.

\subsection*{Formal research question}

For transition kernels compatible with physical evidence and stated economic propagation restrictions, characterize optimal mitigation and tipping-risk values. The following is a conditional robust benchmark within this research family. Suppose $s_t=(x_t,J_t)$ is observed before $a_t$, is sufficient for controlled transitions, and, together with $t$ and $a_t$, determines the conditional ambiguity set $\mathcal Q_t(x_t,J_t,a_t)$. The policymaker observes state and action histories, not nature's selected kernels. Let $h_t=(s_0,a_0,\ldots,a_{t-1},s_t)$ and let $\Pi$ consist of feasible measurable policies $a_t=\pi_t(h_t)$. A nature strategy $\nu$ measurably selects $Q_t^\nu(\cdot\mid h_t,a_t)\in\mathcal Q_t(x_t,J_t,a_t)$ after observing $h_t$ and the chosen action $a_t$. Let $\mathfrak N$ contain all such selections, independently at every date and history, with no additional cross-history or cross-date constraint; this is the pasting convention meant by rectangularity here. Admissible inputs for this benchmark must permit these selections and the path laws $P_s^{\pi,\nu}$ induced from $s_0=s$ by successive conditional kernels. Its welfare objective is
\[
V_0(s)=\sup_{\pi\in\Pi}\inf_{\nu\in\mathfrak N}
\mathbb E_{P_s^{\pi,\nu}}\!\left[
\sum_{t=0}^{H-1}\beta^t U(x_t,J_t,a_t)+\beta^H G(x_H,J_H)
\right].
\]
Under conditions justifying dynamic programming and the required measurable selections, its time-consistent Bellman recursion is
\[
V_t(x,j)=\sup_{a\in A(x,j)}
\left\{U(x,j,a)+\beta\inf_{Q\in\mathcal Q_t(x,j,a)}
\mathbb E_Q[V_{t+1}(x',j')]\right\},\quad t=0,\ldots,H-1,\quad V_H(x,j)=G(x,j).
\]
Here $\mathcal Q_t$ is an explicitly specified nonempty set of transition kernels satisfying the support restriction above. Full observation and state-dependent ambiguity are conditions for this displayed benchmark, not consequences of history rectangularity alone. If the conditional sets depend on more of the history, use that history or a justified sufficient information state; if observation is partial, supply the observation process and admissible information-based policies before asserting a Bellman representation.

Establish conditions for finite well-defined values and measurable $\varepsilon$-optimal policies for every $\varepsilon>0$; distinguish these conclusions from attainment by optimal policies or worst-case kernels. State the domain, feasibility, integrability, regularity and measurable-selection assumptions of each result. For each approximation result, specify the approximation scheme, error metric, tolerance and whether the bound is pointwise or uniform over a declared primitive class. Distinguish value error, policy welfare loss and risk error; for sampling-based bounds also state the probability and confidence criterion. Report dependence on horizon, discounting, joint-kernel approximation and other relevant primitive or approximation parameters. These conditions and bounds are research outputs rather than a requirement to use a finite-state algorithm.

For each tipping-risk result, declare the risk functional, including the event or transition-time variable, horizon, initial condition, policy and treatment of ambiguity. For comparisons of hazard, damage or interaction restrictions, state the baseline and alternative primitive classes, the quantifiers meant by a robust action change, and whether optimal or $\varepsilon$-optimal action sets or a declared selector are compared. Worst-case optimization, agreement across physical models and stability under perturbations are distinct comparisons. When adding migration or capital destruction, supply their transition and feasibility mechanisms, the common baseline and the population and welfare-weight conventions needed to compare the resulting policy conclusions.

For monetization, declare the reference consumption allocation, transfer date and incidence, welfare weights, price basis and real-currency year. Distinguish a local money-metric derivative, which requires differentiability and a positive marginal welfare value for the declared transfer direction, from a finite compensating or equivalent transfer, which requires a stated transfer experiment and welfare indifference equation, with existence and uniqueness addressed. A finite welfare difference divided by a marginal consumption value is only a reference marginal conversion unless an exact finite-change interpretation is separately justified.

\subsection*{Interpretation and scope}

A threshold in a physical state is not automatically a discontinuity in aggregate monetary damages. Economic effects depend on exposure, adaptation and substitution, and a probability of tipping differs from the distribution of transition times. Rectangularity permits nature's conditional choices to vary by history; a single fixed but unknown physical model defines a different ambiguity problem and can invalidate this Bellman recursion. Bayesian learning requires a prior and observation process. A finite-horizon irreversibility approximation must not be presented as a universal physical fact. Sensitivity to discounting and regional welfare weights should be reported separately from estimated hazard uncertainty.

\subsection*{Source audit and status}

[S1] identifies physical tipping elements; [S2] studies optimal policy in a specified tipping IAM. The input's Cai--Judd--Lontzek (2016) does not uniquely match the relevant publication: the verified replacement [S3] is Cai--Lenton--Lontzek (2016), with Judd acknowledged rather than listed as author. This correction is explicit, and intended-source identity remains uncertain. The coupled robust-control target is an editorial research family, not a conjecture asserted by those sources or certified open in 2026.

\subsection*{References}

\begin{enumerate}[label={[\textnormal{S}\arabic*]},leftmargin=*]

\item Timothy M. Lenton, Hermann Held, Elmar Kriegler, Jim W. Hall, Wolfgang Lucht, Stefan Rahmstorf and Hans Joachim Schellnhuber (2008). \emph{Tipping elements in the Earth’s climate system}. Proceedings of the National Academy of Sciences 105(6), 1786–1793. \url{https://doi.org/10.1073/pnas.0705414105}. \textbf{Locator:} Publisher or institutional abstract and bibliographic record. \textbf{Support and limits:} Physical tipping elements and expert assessment; does not identify an economic welfare distribution.

\item Derek Lemoine and Christian Traeger (2014). \emph{Watch Your Step: Optimal Policy in a Tipping Climate}. American Economic Journal: Economic Policy 6(1), 137–166. \url{https://doi.org/10.1257/pol.6.1.137}. \textbf{Locator:} Publisher or institutional abstract and bibliographic record. \textbf{Support and limits:} Specified stochastic IAM with tipping; optimal-tax claims are model conditional.

\item Yongyang Cai, Timothy M. Lenton and Thomas S. Lontzek (2016). \emph{Risk of multiple interacting tipping points should encourage rapid CO2 emission reduction}. Nature Climate Change 6, 520–525. \url{https://doi.org/10.1038/nclimate2964}. \textbf{Locator:} Publisher citation and author list, lines 206–264. \textbf{Support and limits:} Verified relevant replacement has Lenton, not Judd. Original author-year combination remains ambiguous; replacement is disclosed.

\end{enumerate}

\subsection*{Common conventions: Source mathematical conventions}

Unless an entry states otherwise, agent and firm sets are finite. State and action spaces are measurable, strategies and outcome maps are measurable, probability laws are specified, and expectations exist for the targets under discussion. Infinite-horizon discount factors lie in $(0,1)$ and discounted objectives are finite. Norms, tolerances and complexity input sizes must be specified for computational comparisons. Constants or primitives that appear locally are inputs, not empirical facts asserted by this catalogue.

\textbf{Identification.} A statistical model $m\in\mathcal M$ implies an observable law $P_m$ and target $\theta(m)$. At law $P$, its identified set is
\[
\Theta(P;\mathcal M)=\{\theta(m):m\in\mathcal M,\ P_m=P\}.
\]
Point identification holds when this set is a singleton, or equivalently when $P_{m_1}=P_{m_2}$ implies $\theta(m_1)=\theta(m_2)$ for all compatible models. Sharp bounds mean bounds attained, or approached when the boundary is not attained, by compatible models. Writing this set is a definition, not a solution: the substantive task is to establish informative restrictions and derive or estimate the set. An unrestricted-confounding model may give only trivial bounds.

\textbf{Inference.} Identification, estimability and valid uncertainty quantification are separate questions. Fix $\alpha\in(0,1)$. A confidence procedure $C_n$ over a declared law class $\mathcal P$ has asymptotic uniform coverage at level $1-\alpha$ when
\[
\liminf_{n\to\infty}\inf_{P\in\mathcal P}\Pr_P\{\theta(P)\in C_n\}\geq1-\alpha.
\]
For set-identified targets the coverage event must be defined separately, for example coverage of the full identified set. If an entry uses identification-indexed models instead of uniquely indexed laws, the same coverage convention is applied over those models. Pointwise limits do not establish uniform validity, and asymptotic validity does not guarantee finite-sample precision. Sample sizes, dependence structures and weak-identification sequences are part of each application.

\textbf{Counterfactuals.} $Y_i(a)$ is a potential outcome under a specified intervention, including its timing, implementation and versions. It is not $Y_i$ conditional on voluntarily choosing $a$. A full intervention vector permits interference; shorthand scalar treatment notation does not silently impose no interference. Assignment, selection, attrition, anticipation and measurement must be specified. A claim about an instrument's local effect is not automatically a population policy effect.

\textbf{Economic equilibrium.} A counterfactual model includes best responses, constraints, feasibility, market clearing, state transitions and expectations consistent with the stated information. When several equilibria exist, either a declared selector is an input or the target is set-valued. Comparisons across policy regimes must use the stated selection convention. Reduced-form relationships are not presumed invariant after an equilibrium-changing intervention.

\textbf{Optimization and welfare.} Optimization is over the stated feasible set. If attainment is not established, $\sup$ or $\inf$ and $\varepsilon$-optimal choices replace an asserted maximizer or minimizer. For example, with value $V=\sup_{a\in\mathcal A}W(a)<\infty$, an $\varepsilon$-optimal set is $\{a\in\mathcal A:W(a)\geq V-\varepsilon\}$ for $\varepsilon>0$. Benefits, costs, risk and health or environmental outcomes can be aggregated only using declared units and conversion weights. Interpersonal utility weights, the treatment of fiscal transfers and foreign profits, discounting, and the behavioral welfare criterion are explicit inputs. A comparative-static derivative additionally requires existence and appropriate differentiability of the selected counterfactual.

\textbf{Sources.} Local labels [S1], [S2], and so forth restart in every question. Locators identify the relevant abstract, model, section or result at the level actually checked; a bibliographic or abstract-level verification is not represented as inspection of an entire proof. Working papers are distinguished from later publications and changing versions. Source summaries are paraphrased. Source checks are limited to their recorded locators and dates; no common audit date is specified for this entry.
% END ECONOMICS STATEMENT
\end{document}
